\documentclass[10pt,a4paper]{amsart}
\usepackage[margin=19mm]{geometry}
\usepackage{amsmath,amssymb,mathtools}
\usepackage[scr=rsfs,cal=euler]{mathalfa}
\usepackage{xfrac}
\usepackage[unicode,hidelinks,bookmarks=false]{hyperref}
\newcommand{\ie}{i.e.\ }
\newcommand{\cf}{cf.\ }
\newcommand{\resp}{respectively\ }
\newcommand{\Subsection}{\subsection}
\providecommand{\nobreakdash}{\nobreak\hbox{-}}
\setlength{\emergencystretch}{2em}
\multlinegap0pt
\usepackage[shortlabels]{enumitem}
\usepackage{amssymb}
\usepackage{mathtools}
\usepackage{xcolor}
\newtheorem{lemma}{Lemma}
\newcommand{\EE}{\mathbb E}
\newcommand{\PP}{\mathbb P}
\title{The toughness of random graphs}
\author{Guang Li, Wenqian Zhang}
\date{Source: arXiv:2608.31056v2 (CC BY 4.0)}
\begin{document}
\maketitle

\noindent\emph{Source and licence.} The toughness of random graphs. Version arXiv:2608.31056v2 (CC BY 4.0), distributed by arXiv under the Creative Commons Attribution 4.0 International licence, http://creativecommons.org/licenses/by/4.0/, as recorded in the arXiv licence metadata for this exact version and shown on the abstract page https://arxiv.org/abs/2608.31056v2. Full licence notice: \texttt{licenses/arxiv-2608.31056-CC-BY-4.0.txt}.\par\smallskip

\noindent\textit{Editorial scope.} Counted unit: the article's lemma eq:binomial-coefficient-upper (source0.tex lines 143--162). The article's complete original proof of it is delivered with it (source0.tex lines 164--253, one \texttt{proof} environment of 90 lines). No mathematics is written, repaired or re-derived: the statement and the proof are the article's own lines, in the article's own notation, and no retained line is substituted or re-flowed.  No further source-local context is needed: every cross-reference of every retained line resolves inside the two delivered spans.  Nothing was omitted from any retained span, so the package carries no omission record.  No retained line contains a citation, so the wrapper carries no bibliography and no bibliography entry.  No retained line contains a figure environment, an image-inclusion command or drawing code, so no image is delivered and the package declares no asset dependency.  Editorial formatting only, in addition: an AMS \texttt{amsart} preamble that restates the article's own declarations and macros used by the retained lines, namely the environments \texttt{lemma} on separate counters, together with the standard packages those lines need.  The retained environments print in this document as Lemma 1 in order of appearance, whereas the article prints the same items with the numbering of its own full document.  The complete native source of the article as distributed in the arXiv e-print for this exact version is bundled unchanged as \texttt{sources/208982/source0.tex}.
\begin{lemma}\label[lemma]{lem:elementary-estimates}
Let $1\leq r \leq N$ and $X\sim\operatorname{Bin}(N,s)$. Then the following conclusions hold.
\begin{enumerate}[label=\textup{(\roman*)}]
 \item 
 \begin{equation}\label{eq:binomial-coefficient-upper}
  \binom Nr\le\left(\frac{eN}{r}\right)^r.
 \end{equation}
 \item For any integer
 $z\ge\max\{1,Ns\}$,
 \begin{equation}\label{eq:binomial-tail-self-contained}
  \PP(X\ge z)\le\left(\frac{eNs}{z}\right)^z.
 \end{equation}
 \item If $r/N\to0$ (for $N\to\infty$), then
 \begin{equation}\label{eq:uniform-stirling-binomial}
  \log\binom Nr
  =r\log\frac Nr+r
   +O\!\left(\frac{r^2}{N}+\log(r+1)\right).
 \end{equation}
\end{enumerate}
\end{lemma}

\begin{proof}
For (i), using
\[
 \log(r!)=\sum_{i=1}^r\log i
 \ge\int_1^r\log x\,dx
 =r\log r-r+1,
\]
we obtain $r!\ge e(r/e)^r\ge(r/e)^r$.  Since
\[
 \binom Nr
 =\frac{N(N-1)\cdots(N-r+1)}{r!}
 \le\frac{N^r}{r!},
\]
we see
\[
 \binom Nr\le\frac{N^r}{(r/e)^r}
 =\left(\frac{eN}{r}\right)^r,
\]
establishing \eqref{eq:binomial-coefficient-upper}.

For (ii), if $z>N$, the event $X\ge z$ is impossible and the inequality
is immediate. We may therefore assume $1\le z\le N$.
Set $Y=\binom Xz$ as $X$ is integer-valued. Clearly, the event $X\ge z$ is equivalent to the event $Y\ge1$.  Using Markov's
inequality we obtain
\begin{equation}\label{eq:factorial-moment-tail}
 \PP(X\ge z)=\PP(Y\ge 1)\le\EE Y
 =\EE\binom Xz.
\end{equation}
Denote $X=\sum_{i=1}^N\xi_i$, where the $\xi_i$ are independent Bernoulli
variables with success probability $s$.  Since the random variable $\binom Xz$
counts the $z$-subsets of successful trials, we have
\[
 \binom Xz
 =\sum_{J\in\binom{[N]}z}\prod_{i\in J}\xi_i.
\]
Then
\[
 \EE\binom Xz
 =\sum_{J\in\binom{[N]}z}s^z
 =\binom Nzs^z.
\]
Using \eqref{eq:factorial-moment-tail} and (i) with $r=z$, we obtain 
\[
 \PP(X\ge z)
 \le\binom Nzs^z
 \le\left(\frac{eN}{z}\right)^zs^z
 =\left(\frac{eNs}{z}\right)^z,
\]
establishing (ii).  (In fact, the calculation before the final numerical
comparison is valid for every integer $1\le z\le N$; the stated assumption
$z\ge Ns$ is the range in which the displayed bound is useful.)

For (iii), we can assume $r/N\le1/2$ as $r/N\to0$.  It is easy to check that
\begin{equation}\label{eq:binomial-product-expansion}
 \log\binom Nr
 =r\log N-\log(r!)
  +\sum_{i=0}^{r-1}\log\left(1-\frac{i}{N}\right).
\end{equation}
Since $\log x$ is increasing for $x>0$, we have
\[
 \int_1^r\log x\,dx
 \le\sum_{i=2}^r\log i
 \le\int_1^r\log x\,dx+\log r.
\]
It is easy to see 
$$\int_1^r\log x\,dx=r\log r-r+1.$$
Thus, the two-sided estimate above is
\[
 \log(r!)=r\log r-r+O(\log(r+1)),
\]
where $\log(r+1)$ also absorbs the constant term when $r=1$.
For $0\le x\le1/2$, using Taylor's formula we have 
$\log(1-x)=-x+O(x^2)$.  
Hence
\begin{align*}
 \sum_{i=0}^{r-1}\log\left(1-\frac{i}{N}\right)
 &=-\frac1N\sum_{i=0}^{r-1}i
   +O\!\left(\frac1{N^2}\sum_{i=0}^{r-1}i^2\right)\\
 &=-\frac{r(r-1)}{2N}+O\!\left(\frac{r^3}{N^2}\right)
 =O\!\left(\frac{r^2}{N}\right),
\end{align*}
where the last equality uses $r/N\le1/2$.  Inserting the last two estimates
into \eqref{eq:binomial-product-expansion} gives
\[
 \log\binom Nr
 =r\log N-r\log r+r
  +O\!\left(\frac{r^2}{N}+\log(r+1)\right),
\]
establishing (iii).
\end{proof}

\end{document}
